\documentclass[11pt]{article}
\usepackage[utf8]{inputenc}
\usepackage[T1]{fontenc}
\usepackage[italian]{babel}
\usepackage{geometry}
\geometry{a4paper, margin=1in}
\usepackage{xcolor}
\usepackage{listings}
\usepackage{hyperref}

\usepackage{tcolorbox}
\tcbuselibrary{breakable}
\begin{document}

\begin{tcolorbox}[colback=blue!5!white,colframe=blue!75!black,title=\textbf{DOMANDA}, breakable]
Esercizio 4\\
Se si dispone di due PC in grado di dialogare sulla stessa rete IP (oppure un PC per il server e uno smartphone per il client sempre in grado di dialogare tra loro a livello IP) provare ad accedere al server che è su Linux mediante diversi tipi di browser e di sistemi operativi. Cosa si può notare?
\end{tcolorbox}

\begin{tcolorbox}[colback=green!5!white,colframe=green!75!black,title=\textbf{RISPOSTA}, breakable]
\textbf{Esercizio 37 - Test WebSocket da diversi dispositivi}

Effettuando il test di accesso da vari dispositivi (come PC Windows, macOS, smartphone Android o iOS) e browser differenti (Chrome, Firefox, Safari, Edge, ecc.), si possono notare due aspetti principali:

\begin{enumerate}
    \item \textbf{Interoperabilità e Standardizzazione:} L'applicazione funziona in modo identico e fluido su tutti i dispositivi e browser moderni senza necessità di installare software aggiuntivo o applicazioni dedicate sul client. Questo perché le WebSocket sono uno standard web nativamente supportato da tutti i principali browser e operano al di sopra del protocollo \texttt{HTTP}/\texttt{TCP}.
    \item \textbf{Sincronizzazione in tempo reale (Real-time):} Qualsiasi messaggio inviato da uno dei dispositivi (ad esempio dallo smartphone) compare istantaneamente sugli schermi di tutti gli altri dispositivi collegati (ad esempio i PC), confermando che il server riesce a mantenere aperte le connessioni persistenti (full-duplex) con tutti i client simultaneamente e a fare broadcast degli eventi in tempo reale, indipendentemente dal sistema operativo o dal browser utilizzato.
\end{enumerate}

\textbf{Nota operativa:} \\
Perché l'accesso abbia successo, è fondamentale che i client nella rete locale conoscano l'indirizzo IP del server (es. \texttt{192.168.x.x}) e che l'eventuale firewall sul server Linux (es. \texttt{ufw} o \texttt{iptables}) sia configurato per accettare connessioni in entrata sulla porta utilizzata dal server Node.js (es. porta 4000). Altrimenti la pagina non si caricherà.
\end{tcolorbox}

\end{document}
